\chapter{El conjunt de parts i la pertinença interna}\label{cap:parts-interna}

\begin{objectius}
\apartat{Prerequisits} Del capítol~\ref{cap:exponencials}: exponencials, avaluació, transposició i la definició \(P(A):=\mathbf2^A\), amb els punts de \(P(A)\) com a parts d'\(A\). Del capítol~\ref{cap:valors}: les operacions booleanes sobre \(\mathbf2\) i sobre les parts, i el complement. Del capítol~\ref{cap:imatges}: la imatge directa, la seva adjunció amb la imatge inversa, l'estabilitat dels epimorfismes i el lema de pertinença després d'un recobriment.
\apartat{Objectius} En acabar el capítol, el lector hauria de poder:
\begin{itemize}
  \item distingir la part \(m\), el seu nom \(\ulcorner m\urcorner\colon\mathbf1\to P(A)\) i el conjunt \(P(A)\), i construir la pertinença interna \(\in_A\) com la part classificada per l'avaluació;
  \item identificar les aplicacions \(X\to P(A)\) amb les parts de \(X\times A\), i llegir un element generalitzat de \(P(A)\) com una família de parts;
  \item definir les operacions booleanes internes per transposició, calcular-les sobre noms de parts i saber per a quines d'elles la lectura puntual val a tots els estadis;
  \item construir la imatge inversa interna \(P(f)\) i els quantificadors \(\exists_f\) i \(\forall_f\) com a aplicacions, i demostrar que \(\exists_f\dashv P(f)\dashv\forall_f\);
  \item demostrar que la imatge directa commuta amb el canvi de base, i usar-ho per llegir \(\exists_f\) sobre elements generalitzats;
  \item definir la inclusió interna sobre \(P(A)\) amb el quantificador universal i demostrar que és un ordre amb \(\cap\) i \(\cup\) com a ínfim i suprem.
\end{itemize}
\end{objectius}

El capítol anterior va definir el \emph{conjunt de parts} d'un conjunt
\(A\) com l'exponencial
\[
P(A):=\mathbf2^{A}
\]
(definició~\ref{def:conjunt-parts}), i en va identificar els punts globals
amb les aplicacions característiques \(A\to\mathbf2\), és a dir, amb les
parts de \(A\) llevat d'equivalència. Amb això, el conjunt de parts va
passar a existir com a \emph{objecte} de la teoria. Però definir un
objecte no és el mateix que estudiar-lo: quedaven obertes la pertinença
interna, les operacions booleanes sobre \(P(A)\) i l'ordre d'inclusió
(observació~\ref{obs:parts-obert}). Aquest capítol tanca aquestes
qüestions.

No hi introduirem cap axioma. Tot el que segueix és conseqüència de
l'existència d'exponencials del capítol~\ref{cap:exponencials}, del
classificador de parts del capítol~\ref{cap:parts} i de les operacions
booleanes sobre \(\mathbf2\) del capítol~\ref{cap:valors}. És, doncs, un
capítol de consolidació: fa interns, com a estructura sobre un únic
conjunt \(P(A)\), els fenòmens que fins ara vivien dispersos entre el
metallenguatge i el conjunt de valors de veritat.

Cal des del principi mantenir separades tres coses que la lectura clàssica
tendeix a confondre. Primer, una \emph{part} és una fletxa
\(m\colon S\hookrightarrow A\), un monomorfisme; és una dada del
metallenguatge. Segon, aquesta part té un únic \emph{punt classificador}
\[
\ulcorner m\urcorner\colon\mathbf1\to P(A),
\]
el nom intern de la seva aplicació característica \(\chi_m\colon A\to\mathbf2\);
és un element global d'un conjunt. Tercer, \(P(A)\) és el \emph{conjunt}
que reuneix tots aquests noms. La distinció entre la part, el seu nom i el
conjunt que els conté és el fil conductor del capítol.

\section{La relació de pertinença interna}\label{sec:pertinenca-interna}

L'avaluació de l'exponencial \(\mathbf2^{A}\) és una aplicació
\[
\operatorname{ev}\colon P(A)\times A\to\mathbf2,
\]
és a dir, una aplicació cap al conjunt de valors de veritat. Pel
classificador de parts (axioma~\ref{ax:representacio}), tota aplicació cap
a \(\mathbf2\) és la característica d'una part del seu domini. La part que
classifica l'avaluació és la relació de pertinença que buscàvem.

\begin{definicio}[Pertinença interna]\index{Pertinença interna}\label{def:pertinenca-interna}
Sigui \(A\) un conjunt. La \emph{relació de pertinença interna} de \(A\)
és la part de \(A\times P(A)\)
\[
{\in_A}\;\colon\;{\in_A}\hookrightarrow A\times P(A)
\]
classificada per l'avaluació, llegida sobre \(A\times P(A)\) via
l'isomorfisme de commutació \(\sigma\colon P(A)\times A\to A\times P(A)\).
És a dir, \(\in_A\) és la part de \(A\times P(A)\) la característica de la
qual és
\[
\chi_{\in_A}=\operatorname{ev}\circ\sigma^{-1}\colon A\times P(A)\to\mathbf2 .
\]
\end{definicio}

\begin{notacio}[Pertinença interna]\label{not:pertinenca-interna}
Escriurem \(\in_A\) tant per a la part com per al seu domini, seguint la
convenció de la relació inversa i la diagonal del
capítol~\ref{cap:relacions}. El predicat eliminable corresponent és
\(\operatorname{Car}(\operatorname{ev}\circ\sigma^{-1},{\in_A},A\times
P(A))\) de la definició~\ref{def:caracteristica}, amb \(\sigma\)
introduïda pel predicat de producte \(\operatorname{Prod}\). Per a un
parell d'elements generalitzats \(a\colon X\to A\) i \(\alpha\colon X\to
P(A)\), escriurem
\[
a\mathrel{\varepsilon}\alpha
\quad\text{en lloc de}\quad
\langle a,\alpha\rangle\in_X{\in_A},
\]
i llegirem «\(a\) pertany internament a \(\alpha\)».
\end{notacio}

La definició no fa res més que aplicar el classificador a una aplicació ja
existent. El seu contingut apareix quan la traduïm a la pertinença
generalitzada del capítol~\ref{cap:parts}.

\begin{proposicio}[Lectura de la pertinença interna]\label{prop:pertinenca-interna-lectura}
Siguin \(a\colon X\to A\) i \(\alpha\colon X\to P(A)\) dos elements
generalitzats a un mateix estadi \(X\). Aleshores
\[
a\mathrel{\varepsilon}\alpha
\quad\Longleftrightarrow\quad
\operatorname{ev}\circ\langle\alpha,a\rangle=1\circ{!}_X .
\]
En particular, si \(\alpha=\ulcorner m\urcorner\circ{!}_X\) és constant amb
valor el nom d'una part \(m\colon S\hookrightarrow A\), aleshores
\[
a\mathrel{\varepsilon}\alpha
\quad\Longleftrightarrow\quad
a\in_X m .
\]
\end{proposicio}

\begin{prova}
Per la propietat característica (notació~\ref{not:caracteristica}),
\(\langle a,\alpha\rangle\in_X{\in_A}\) equival a
\(\chi_{\in_A}\circ\langle a,\alpha\rangle=1\circ{!}_X\). Ara bé,
\(\chi_{\in_A}=\operatorname{ev}\circ\sigma^{-1}\) i
\(\sigma^{-1}\circ\langle a,\alpha\rangle=\langle\alpha,a\rangle\), d'on
la primera equivalència.

Per a la segona, sigui \(\alpha=\ulcorner m\urcorner\circ{!}_X\). El punt
classificador \(\ulcorner m\urcorner\) és, per definició, el nom intern de
\(\chi_m\), de manera que descurrificar-lo recupera \(\chi_m\): la
transposada de \(\operatorname{ev}\circ(\ulcorner m\urcorner\times\id_A)\)
és \(\ulcorner m\urcorner\), i per la unicitat de la transposada
(definició~\ref{def:exponencial}) es té
\(\operatorname{ev}\circ(\ulcorner m\urcorner\times\id_A)=\chi_m\circ\rho\),
on \(\rho\colon\mathbf1\times A\to A\) és la projecció. Component amb
\(\langle{!}_X,a\rangle\) i usant \(\alpha=\ulcorner m\urcorner\circ{!}_X\),
\[
\operatorname{ev}\circ\langle\alpha,a\rangle
=\chi_m\circ a .
\]
Així \(a\mathrel{\varepsilon}\alpha\) equival a
\(\chi_m\circ a=1\circ{!}_X\), que és exactament \(a\in_X m\).
\end{prova}

\begin{observacio}[La part, el seu nom i el conjunt]\label{obs:tres-nivells}
La proposició~\ref{prop:pertinenca-interna-lectura} és el lloc on es
paguen els dividends de mantenir separats els tres nivells. La
\emph{part} \(m\colon S\hookrightarrow A\) és una fletxa; el seu
\emph{nom} \(\ulcorner m\urcorner\colon\mathbf1\to P(A)\) és un punt del
conjunt \(P(A)\); i la pertinença \(a\in_X m\), que abans era una
factorització a través de \(m\), es converteix ara en la pertinença
interna \(a\mathrel{\varepsilon}\ulcorner m\urcorner\), una condició
expressada íntegrament amb aplicacions cap a \(P(A)\) i \(\mathbf2\). La
relació primitiva \(\in\) de la teoria clàssica, que allà és una noció no
definida entre individus, apareix aquí com una \emph{part construïda},
\(\in_A\hookrightarrow A\times P(A)\), classificada per l'avaluació.
\end{observacio}

\begin{observacio}[Lectura clàssica]\label{obs:pertinenca-interna-classica}
En conjunts, \(P(A)\) és el conjunt de subconjunts de \(A\), un punt
\(\alpha\in P(A)\) és un subconjunt concret, i
\(\operatorname{ev}(\alpha,a)\) val \(1\) exactament quan \(a\in\alpha\).
La part \(\in_A\subseteq A\times P(A)\) és, doncs, el graf de la relació
de pertinença: el conjunt de parells \((a,\alpha)\) amb \(a\in\alpha\). El
que en la presentació habitual és l'axioma que dona sentit a l'expressió
\(x\in y\), aquí és un teorema sobre una part determinada per l'avaluació.
\end{observacio}

\section{\texorpdfstring{Aplicacions cap a \(P(A)\) i parts de \(X\times A\)}{Aplicacions cap a P(A) i parts de X per A}}\label{sec:parts-parametritzades}

El teorema dels punts de l'exponencial (teorema~\ref{teo:punts-exponencial})
identificava els punts globals \(\mathbf1\to P(A)\) amb les parts de
\(A\). El cas general d'un estadi \(X\) qualsevol dona la versió
parametritzada: una aplicació \(X\to P(A)\) és una \emph{família de parts
de \(A\) indexada per \(X\)}, o equivalentment una part de \(X\times A\).

\begin{teorema}[Parts parametritzades]\label{teo:parts-parametritzades}
Per a conjunts \(X\) i \(A\), la transposició estableix una
correspondència biunívoca, natural en \(X\), entre:
\begin{enumerate}
\item les aplicacions \(\alpha\colon X\to P(A)\);
\item les aplicacions \(X\times A\to\mathbf2\);
\item les parts de \(X\times A\) llevat d'equivalència.
\end{enumerate}
La correspondència entre \((1)\) i \((2)\) és la descurrificació; la
correspondència entre \((2)\) i \((3)\) és el classificador de parts.
\end{teorema}

\begin{prova}
L'equivalència entre \((1)\) i \((2)\) és la propietat universal de
l'exponencial \(P(A)=\mathbf2^{A}\)
(definició~\ref{def:exponencial}) aplicada amb codomini \(\mathbf2\): a
cada \(\alpha\colon X\to P(A)\) li correspon
\(\operatorname{ev}\circ(\alpha\times\id_A)\colon X\times A\to\mathbf2\), i
la transposició és biunívoca per la unicitat de \(\widehat h\). La
naturalitat en \(X\) és la de la proposició~\ref{prop:families},
particularitzada a \(B=\mathbf2\).

L'equivalència entre \((2)\) i \((3)\) és l'axioma de representació
(axioma~\ref{ax:representacio}): tota aplicació \(X\times A\to\mathbf2\) és
la característica d'una única part de \(X\times A\) llevat d'equivalència,
i tota part de \(X\times A\) té una única característica. Componen les dues
correspondències.
\end{prova}

\begin{observacio}[Elements generalitzats de \(P(A)\)]\label{obs:generalitzats-parts}
El teorema diu que un element generalitzat de \(P(A)\) a l'estadi \(X\)
---una aplicació \(X\to P(A)\)--- no és una sola part de \(A\), sinó una
part de \(X\times A\), és a dir, una família de parts de \(A\)
parametritzada per \(X\). El cas global \(X=\mathbf1\) recupera una única
part de \(\mathbf1\times A\cong A\)
(teorema~\ref{teo:punts-exponencial}). Aquesta és la mateixa disciplina
que regeix tot el desenvolupament: les nocions es defineixen sobre
elements generalitzats, i el cas global n'és el particular \(X=\mathbf1\).
\end{observacio}

\begin{corollari}[La pertinença interna és universal]\label{cor:pertinenca-universal}
Sigui \(\alpha\colon X\to P(A)\) i sigui \(s_\alpha\hookrightarrow
X\times A\) la part que li correspon pel
teorema~\ref{teo:parts-parametritzades}. Aleshores, per a tot estadi
\(Y\) i tot parell \(a\colon Y\to A\), \(x\colon Y\to X\),
\[
a\mathrel{\varepsilon}(\alpha\circ x)
\quad\Longleftrightarrow\quad
\langle x,a\rangle\in_Y s_\alpha .
\]
En particular, la part \(\in_A\hookrightarrow A\times P(A)\) és la part
associada a la identitat \(\id_{P(A)}\colon P(A)\to P(A)\).
\end{corollari}

\begin{prova}
Per la proposició~\ref{prop:pertinenca-interna-lectura},
\(a\mathrel{\varepsilon}(\alpha\circ x)\) equival a
\(\operatorname{ev}\circ\langle\alpha\circ x,a\rangle=1\circ{!}_Y\). Ara
\(\langle\alpha\circ x,a\rangle=(\alpha\times\id_A)\circ\langle x,a\rangle\),
de manera que la condició és
\(\operatorname{ev}\circ(\alpha\times\id_A)\circ\langle x,a\rangle
=1\circ{!}_Y\); i \(\operatorname{ev}\circ(\alpha\times\id_A)\) és, per
construcció, la característica de \(s_\alpha\). Això és
\(\langle x,a\rangle\in_Y s_\alpha\). Prenent \(X=P(A)\) i
\(\alpha=\id_{P(A)}\), la part associada té característica
\(\operatorname{ev}\), que sobre \(A\times P(A)\) és la de \(\in_A\).
\end{prova}

\begin{observacio}[Lectura clàssica]\label{obs:parts-parametritzades-classica}
Les tres presentacions són les que un curs elemental usa
indistintament: una funció \(\alpha\colon X\to\mathcal P(A)\) que a cada
índex assigna un subconjunt; una relació \(R\subseteq X\times A\); i una
funció \((x,a)\mapsto\chi_R(x,a)\) amb valors de veritat. Que siguin la
mateixa dada és el que legitima escriure indistintament \(a\in\alpha_x\),
\((x,a)\in R\) o \(\chi_R(x,a)=1\). El teorema
dona a aquesta equivalència, habitualment implícita, un estatut
estructural.
\end{observacio}

\section{\texorpdfstring{Operacions booleanes internes sobre \(P(A)\)}{Operacions booleanes internes sobre P(A)}}\label{sec:booleanes-internes}

Al capítol~\ref{cap:valors} les operacions booleanes es van definir com a
aplicacions sobre el conjunt de valors de veritat: la conjunció i la
disjunció \(\wedge,\vee\colon\mathbf2\times\mathbf2\to\mathbf2\)
(definició~\ref{def:conjuncio-disjuncio}) i la negació
\(\neg\colon\mathbf2\to\mathbf2\) (definició~\ref{def:negacio}). Aquestes
operacions actuen sobre valors de veritat, un a un. Ara les elevarem a
operacions sobre \(P(A)\), on cada punt reuneix una part sencera. La clau
és que les operacions sobre \(\mathbf2\) actuen «puntualment» en cada
element de \(A\), i la transposició les converteix en operacions sobre el
conjunt de parts.

\begin{definicio}[Operacions booleanes internes]\index{Operacions booleanes internes}\label{def:booleanes-internes}
Sigui \(A\) un conjunt. Definim les aplicacions internes següents sobre
\(P(A)\), com a transposades de les operacions puntuals corresponents:
\begin{itemize}
\item la \emph{intersecció interna}
\(\cap_A\colon P(A)\times P(A)\to P(A)\), transposada de
\[
P(A)\times P(A)\times A
\xrightarrow{\;\langle\operatorname{ev}_1,\operatorname{ev}_2\rangle\;}
\mathbf2\times\mathbf2
\xrightarrow{\;\wedge\;}
\mathbf2 ;
\]
\item la \emph{unió interna}
\(\cup_A\colon P(A)\times P(A)\to P(A)\), transposada de la mateixa
aplicació amb \(\vee\) en lloc de \(\wedge\);
\item el \emph{complement intern}
\(\neg_A\colon P(A)\to P(A)\), transposada de
\(\neg\circ\operatorname{ev}\colon P(A)\times A\to\mathbf2\).
\end{itemize}
Aquí \(\operatorname{ev}_1\) i \(\operatorname{ev}_2\) denoten les
avaluacions de la primera i la segona còpia de \(P(A)\) sobre \(A\),
obtingudes component \(\operatorname{ev}\) amb les projeccions adequades.
Anàlogament, els \emph{punts distingits} \(\varnothing_A,\;A_A\colon
\mathbf1\to P(A)\) són els noms de les aplicacions constants
\(0\circ{!}_A\) i \(1\circ{!}_A\) (proposició~\ref{prop:caracteristica-total}),
és a dir, de la part buida i la part total del capítol~\ref{cap:valors}.
\end{definicio}

\begin{notacio}[Operacions internes]\label{not:booleanes-internes}
Per a elements generalitzats \(\alpha,\beta\colon X\to P(A)\) escriurem
\(\alpha\cap_A\beta\), \(\alpha\cup_A\beta\) i \(\neg_A\alpha\) per a les
composicions de \(\langle\alpha,\beta\rangle\) o \(\alpha\) amb les
aplicacions de la definició~\ref{def:booleanes-internes}. Quan el conjunt
\(A\) sigui inequívoc, ometrem el subíndex.
\end{notacio}

El sentit d'aquestes operacions és que calculen, punt a punt, la
pertinença interna al resultat. Aquesta és la propietat que les
caracteritza.

\begin{proposicio}[Càlcul de les operacions internes]\label{prop:booleanes-puntuals}
Siguin \(\alpha,\beta\colon X\to P(A)\) i \(a\colon X\to A\). Les
operacions internes es calculen component les avaluacions amb les
operacions sobre \(\mathbf2\):
\begin{align*}
\operatorname{ev}\circ\langle\alpha\cap\beta,a\rangle
&=\wedge\circ\langle\operatorname{ev}\circ\langle\alpha,a\rangle,
\operatorname{ev}\circ\langle\beta,a\rangle\rangle,\\
\operatorname{ev}\circ\langle\alpha\cup\beta,a\rangle
&=\vee\circ\langle\operatorname{ev}\circ\langle\alpha,a\rangle,
\operatorname{ev}\circ\langle\beta,a\rangle\rangle,\\
\operatorname{ev}\circ\langle\neg\alpha,a\rangle
&=\neg\circ\operatorname{ev}\circ\langle\alpha,a\rangle .
\end{align*}
En conseqüència, a tot estadi \(X\),
\[
a\mathrel{\varepsilon}(\alpha\cap\beta)
\quad\Longleftrightarrow\quad
a\mathrel{\varepsilon}\alpha\ \text{ i }\ a\mathrel{\varepsilon}\beta ;
\]
i sobre punts globals (\(X=\mathbf1\)), a més,
\(a\mathrel{\varepsilon}(\alpha\cup\beta)\) si i només si
\(a\mathrel{\varepsilon}\alpha\) o \(a\mathrel{\varepsilon}\beta\), i
\(a\mathrel{\varepsilon}(\neg\alpha)\) si i només si no
\(a\mathrel{\varepsilon}\alpha\). Finalment, sobre els noms de dues parts
\(m,n\) de \(A\), les operacions internes recuperen les operacions sobre
parts del capítol~\ref{cap:valors}:
\[
\ulcorner m\urcorner\cap\ulcorner n\urcorner=\ulcorner m\cap n\urcorner,
\qquad
\ulcorner m\urcorner\cup\ulcorner n\urcorner=\ulcorner m\cup n\urcorner,
\qquad
\neg\ulcorner m\urcorner=\ulcorner m^{c}\urcorner .
\]
\end{proposicio}

\begin{prova}
Les tres identitats són la descurrificació de la
definició~\ref{def:booleanes-internes}: per la propietat universal de
l'exponencial,
\(\operatorname{ev}\circ((\alpha\cap\beta)\times\id_A)
=\wedge\circ\langle\operatorname{ev}_1,\operatorname{ev}_2\rangle\circ(\langle\alpha,\beta\rangle\times\id_A)\),
i component amb \(\langle\id_X,a\rangle\) s'obté la primera; les altres
dues són idèntiques.

Per a la lectura de \(\cap\): \(\wedge\) és la característica de
\(\delta_{11}=\langle1,1\rangle\colon\mathbf1\hookrightarrow\mathbf2\times\mathbf2\)
(definició~\ref{def:conjuncio-disjuncio}), de manera que
\(\wedge\circ\langle u,v\rangle=1\circ{!}_X\) si i només si
\(\langle u,v\rangle\) es factoritza per \(\delta_{11}\), és a dir, si i
només si \(u=1\circ{!}_X\) i \(v=1\circ{!}_X\); amb
\(u=\operatorname{ev}\circ\langle\alpha,a\rangle\) i
\(v=\operatorname{ev}\circ\langle\beta,a\rangle\), la
proposició~\ref{prop:pertinenca-interna-lectura} dona l'equivalència, a
qualsevol estadi.

Per a \(\cup\) i \(\neg\) la lectura amb \enquote{o} i \enquote{no} només
val sobre punts globals: allà cada avaluació \(u\) i \(v\) és el punt
\(0\) o el punt \(1\) (corol·lari~\ref{cor:elements-de-2}), i les taules
de veritat (proposició~\ref{prop:taules-veritat}) decideixen. Sobre un estadi
general no val: per exemple, amb \(X=\mathbf2\), \(a\) tal que
\(\operatorname{ev}\circ\langle\alpha,a\rangle=\id_{\mathbf2}\) i
\(\operatorname{ev}\circ\langle\beta,a\rangle=\neg\), es té
\(a\mathrel{\varepsilon}(\alpha\cup\beta)\) sense que
\(a\mathrel{\varepsilon}\alpha\) ni \(a\mathrel{\varepsilon}\beta\).
És la mateixa frontera que separa el complement d'una part de la negació
puntual (observació~\ref{obs:complement-classic}).

Per a la darrera afirmació no cal cap lectura puntual: les parts
\(m\cap n\), \(m\cup n\) i \(m^c\) es van definir per les seves
característiques, \(\wedge\circ\langle\chi_m,\chi_n\rangle\),
\(\vee\circ\langle\chi_m,\chi_n\rangle\) i \(\neg\circ\chi_m\)
(proposició~\ref{prop:caracteristica-interseccio}, definició~\ref{def:operacions-parts}, proposició~\ref{prop:caracteristica-complement}), i les tres identitats
anteriors amb \(X=\mathbf1\), \(\alpha=\ulcorner m\urcorner\) i
\(\beta=\ulcorner n\urcorner\) diuen que la descurrificació de
\(\ulcorner m\urcorner\cap\ulcorner n\urcorner\) és exactament
\(\wedge\circ\langle\chi_m,\chi_n\rangle\), i anàlogament per a les
altres dues. Per la unicitat del nom (teorema~\ref{teo:punts-exponencial}),
els punts coincideixen.
\end{prova}

\begin{observacio}[Lectura clàssica]\label{obs:booleanes-internes-classica}
Les operacions internes són les operacions de conjunts habituals,
\(\alpha\cap\beta\), \(\alpha\cup\beta\) i \(A\setminus\alpha\), vistes ara
com a \emph{aplicacions} \(P(A)\times P(A)\to P(A)\) i \(P(A)\to P(A)\), i
no com a operacions definides element a element de manera externa. El
guany estructural és que l'àlgebra de subconjunts de \(A\), que en el
metallenguatge és una àlgebra booleana sobre les classes d'equivalència de
parts, es realitza ara com una estructura algebraica \emph{sobre un únic
conjunt} \(P(A)\), amb operacions que són aplicacions de la teoria.
\end{observacio}

\begin{observacio}[Les lleis booleanes es transporten]\label{obs:lleis-transportades}
Les lleis booleanes de les parts (teorema~\ref{teo:lleis-booleanes})
---commutativitat, associativitat, distributivitat, lleis de De~Morgan,
complementació--- valen també per a les operacions internes de
\(P(A)\). La raó és que cada llei és una igualtat entre aplicacions cap a
\(\mathbf2\) que ja es compleix puntualment sobre \(\mathbf2\)
(les taules de veritat), i la transposició, en ser biunívoca, transporta
la igualtat a la corresponent entre aplicacions cap a \(P(A)\). Aquest
capítol no reescriu les demostracions: n'hi ha prou de fer notar que la
unicitat de la transposada converteix cada identitat booleana puntual en
la identitat interna corresponent.
\end{observacio}

\section{Imatge inversa, imatge directa i quantificadors}\label{sec:quantificadors}

Una aplicació \(f\colon A\to B\) indueix operacions entre els conjunts de
parts \(P(A)\) i \(P(B)\). La imatge inversa del capítol~\ref{cap:parts} i
la imatge directa del capítol~\ref{cap:imatges}, que allà eren operacions
sobre parts individuals, es realitzen ara com a aplicacions internes; i
l'adjunció entre elles s'estén a una parella de quantificadors associats
a \(f\).

\begin{definicio}[Imatge inversa interna]\index{Imatge inversa interna}\label{def:imatge-inversa-interna}
Sigui \(f\colon A\to B\). La \emph{imatge inversa interna} és l'aplicació
\[
P(f)\colon P(B)\to P(A)
\]
transposada de
\(\operatorname{ev}\circ(\id_{P(B)}\times f)\colon P(B)\times A\to\mathbf2\).
Equivalentment, \(P(f)\) és l'única aplicació tal que, per a tot
\(\beta\colon X\to P(B)\) i tot \(a\colon X\to A\),
\[
a\mathrel{\varepsilon}\bigl(P(f)\circ\beta\bigr)
\quad\Longleftrightarrow\quad
(f\circ a)\mathrel{\varepsilon}\beta .
\]
\end{definicio}

\begin{proposicio}[La imatge inversa interna estén \(f^{-1}\)]\label{prop:imatge-inversa-interna}
Sobre els noms de parts, \(P(f)\) recupera la imatge inversa del
capítol~\ref{cap:parts}: per a tota part \(n\) de \(B\),
\[
P(f)\circ\ulcorner n\urcorner=\ulcorner f^{-1}(n)\urcorner .
\]
Més generalment, si \(\beta\colon X\to P(B)\) correspon a la part
\(s_\beta\) de \(X\times B\) (teorema~\ref{teo:parts-parametritzades}),
aleshores \(P(f)\circ\beta\) correspon a la part
\((\id_X\times f)^{-1}(s_\beta)\) de \(X\times A\).
\end{proposicio}

\begin{prova}
La descurrificació de \(P(f)\circ\beta\) és
\(\operatorname{ev}\circ(\id_{P(B)}\times f)\circ(\beta\times\id_A)
=\operatorname{ev}\circ(\beta\times\id_B)\circ(\id_X\times f)
=\chi_{s_\beta}\circ(\id_X\times f)\),
que és la característica de \((\id_X\times f)^{-1}(s_\beta)\)
(proposició~\ref{prop:caracteristica-imatge-inversa}). El cas dels noms és
\(X=\mathbf1\), amb \(\mathbf1\times A\cong A\): la descurrificació de
\(P(f)\circ\ulcorner n\urcorner\) és \(\chi_n\circ f=\chi_{f^{-1}(n)}\).
\end{prova}

\begin{proposicio}[La imatge inversa commuta amb el complement]\label{prop:imatge-inversa-complement}
Sigui \(f\colon A\to B\). Aleshores \(P(f)\circ\neg_B=\neg_A\circ P(f)\).
\end{proposicio}

\begin{prova}
Comparem descurrificacions, com a aplicacions \(P(B)\times A\to\mathbf2\).
D'una banda,
\(\operatorname{ev}\circ((P(f)\circ\neg_B)\times\id_A)
=\operatorname{ev}\circ(\id\times f)\circ(\neg_B\times\id_A)
=\operatorname{ev}\circ(\neg_B\times\id_B)\circ(\id\times f)
=\neg\circ\operatorname{ev}\circ(\id\times f)\).
De l'altra,
\(\operatorname{ev}\circ((\neg_A\circ P(f))\times\id_A)
=\neg\circ\operatorname{ev}\circ(P(f)\times\id_A)
=\neg\circ\operatorname{ev}\circ(\id\times f)\).
Per la unicitat de la transposada, les dues aplicacions coincideixen.
\end{prova}

Per construir el quantificador existencial com una \emph{aplicació}
\(P(A)\to P(B)\) necessitem la imatge directa del capítol~\ref{cap:imatges}
aplicada no a una part d'\(A\), sinó a la part universal \(\in_A\), i un
lema que diu que la imatge directa es comporta bé sota canvi de base. És
la propietat que fa que la lectura sobre elements generalitzats de
\(\exists_f\) sigui la que s'espera.

\begin{lema}[La imatge directa commuta amb el canvi de base]\label{lema:beck-chevalley}
Sigui
\[
\begin{aligned}
A' &\xrightarrow{\;h\;} A\\
{\scriptstyle f'}\big\downarrow\ \ &\qquad\big\downarrow{\scriptstyle f}\\
B' &\xrightarrow{\;h'\;} B
\end{aligned}
\]
un quadrat cartesià. Per a tota part \(s\colon S\hookrightarrow A\),
\[
(h')^{-1}\bigl(f_{*}s\bigr)\;\equiv_{B'}\;f'_{*}\bigl(h^{-1}(s)\bigr).
\]
\emph{Axiomes usats:} els del capítol~\ref{cap:imatges} (en particular
l'estabilitat dels epimorfismes, que usa la separació global); cap elecció.
\end{lema}

\begin{prova}
Factoritzem \(f\circ s=m\circ e\) amb \(e\colon S\twoheadrightarrow I\)
epimorfisme i \(m\colon I\hookrightarrow B\) la imatge \(f_*s\)
(teorema~\ref{teo:factoritzacio}). Fem el pullback de tota aquesta
factorització al llarg de \(h'\): obtenim \(I':=(h')^{-1}(m)\), una part
de \(B'\), i \(S':=B'\times_B S\) amb una aplicació \(e'\colon S'\to I'\)
que és el pullback de \(e\) al llarg de \(I'\to I\), i per tant un
epimorfisme (lema~\ref{lema:estabilitat-epi-regular}). D'altra banda,
com que el quadrat de l'enunciat és cartesià, \(S'=B'\times_BS\) coincideix
amb \(A'\times_AS\) (dos quadrats cartesians encaixats formen un quadrat
cartesià), que és el domini de \(h^{-1}(s)\); i la composició
\(f'\circ h^{-1}(s)\colon S'\to B'\) és \(I'\hookrightarrow B'\) precedida
de \(e'\). Tenim, doncs, una factorització de \(f'\circ h^{-1}(s)\) com a
epimorfisme seguit de monomorfisme, i per la unicitat de la factorització
imatge, \(f'_*(h^{-1}(s))\equiv I'=(h')^{-1}(f_*s)\).
\end{prova}

\begin{definicio}[Quantificadors associats a una aplicació]\index{Quantificadors associats a una aplicació}\label{def:quantificadors}
Sigui \(f\colon A\to B\). Sigui \(E_A\hookrightarrow P(A)\times A\) la
part classificada per l'avaluació (la pertinença interna amb els factors
intercanviats: \(E_A=\sigma^{-1}\circ{\in_A}\)). El \emph{quantificador
existencial} associat a \(f\) és l'aplicació
\[
\exists_f\colon P(A)\to P(B)
\]
transposada de la característica de la imatge directa
\((\id_{P(A)}\times f)_{*}E_A\hookrightarrow P(A)\times B\). El
\emph{quantificador universal} és
\[
\forall_f:=\neg_B\circ\exists_f\circ\neg_A\colon P(A)\to P(B).
\]
Escriurem \(\exists_A:=\exists_{{!}_A}\) i \(\forall_A:=\forall_{{!}_A}\)
per a l'única aplicació \({!}_A\colon A\to\mathbf1\); són aplicacions
\(P(A)\to P(\mathbf1)\cong\mathbf2\).
\end{definicio}

\begin{proposicio}[Lectura dels quantificadors]\label{prop:quantificadors-lectura}
Sigui \(f\colon A\to B\) i sigui \(\alpha\colon X\to P(A)\), amb part
associada \(s_\alpha\hookrightarrow X\times A\). Aleshores:
\begin{enumerate}
\item \(\exists_f\circ\alpha\) correspon a la part
\((\id_X\times f)_{*}(s_\alpha)\) de \(X\times B\), i
\(\forall_f\circ\alpha\) a la part
\(\bigl((\id_X\times f)_{*}(s_\alpha^{c})\bigr)^{c}\). En particular, sobre
noms, \(\exists_f\ulcorner m\urcorner=\ulcorner f_*m\urcorner\).
\item Per a \(b\colon X\to B\), \(b\mathrel{\varepsilon}(\exists_f\circ\alpha)\)
si i només si existeixen un recobriment epimòrfic
\(p\colon X'\twoheadrightarrow X\) i un \(a\colon X'\to A\) amb
\(f\circ a=b\circ p\) i \(a\mathrel{\varepsilon}(\alpha\circ p)\). Sobre
punts globals, \(b\mathrel{\varepsilon}(\exists_f\circ\alpha)\) si i només
si hi ha un punt \(a\) amb \(f\circ a=b\) i \(a\mathrel{\varepsilon}\alpha\).
\item Per a \(b\colon X\to B\), \(b\mathrel{\varepsilon}(\forall_f\circ\alpha)\)
si i només si, per a tot \(x\colon Y\to X\) i tot \(a\colon Y\to A\) amb
\(f\circ a=b\circ x\), es té \(a\mathrel{\varepsilon}(\alpha\circ x)\).
\end{enumerate}
Amb l'axioma d'elecció, el recobriment de~(2) és escindit i el testimoni
\(a\) es pot prendre ja sobre \(X\).
\end{proposicio}

\begin{prova}
\emph{(1)} La descurrificació de \(\exists_f\circ\alpha\) és
\(\chi_{(\id\times f)_*E_A}\circ(\alpha\times\id_B)\), la característica
de \((\alpha\times\id_B)^{-1}\bigl((\id_{P(A)}\times f)_*E_A\bigr)\). El
quadrat format per \(\alpha\times\id_A\colon X\times A\to P(A)\times A\),
\(\id_X\times f\), \(\id_{P(A)}\times f\) i \(\alpha\times\id_B\) és
cartesià (és el producte del quadrat cartesià trivial de vèrtex \(X\)
sobre \(P(A)\) amb el de \(f\) sobre ell mateix). Pel
lema~\ref{lema:beck-chevalley},
\((\alpha\times\id_B)^{-1}\bigl((\id\times f)_*E_A\bigr)
\equiv(\id_X\times f)_*\bigl((\alpha\times\id_A)^{-1}E_A\bigr)\), i
\((\alpha\times\id_A)^{-1}E_A\) té característica
\(\operatorname{ev}\circ(\alpha\times\id_A)\), que és la de \(s_\alpha\).
Per a \(\forall_f\), apliquem el que acabem de veure a \(\neg_A\circ\alpha\),
que correspon a \(s_\alpha^c\) (proposició~\ref{prop:booleanes-puntuals}),
i després \(\neg_B\). El cas dels noms és \(X=\mathbf1\).

\emph{(2)} Per~(1), \(b\mathrel{\varepsilon}(\exists_f\circ\alpha)\) vol dir
\(\langle\id_X,b\rangle\in_X(\id_X\times f)_*(s_\alpha)\). Si hi ha
\(p\) i \(a\) com a l'enunciat, aleshores
\(\langle p,a\rangle\in_{X'}s_\alpha\) (corol·lari~\ref{cor:pertinenca-universal})
i \((\id_X\times f)\circ\langle p,a\rangle=\langle p,f\circ a\rangle
=\langle\id_X,b\rangle\circ p\), de manera que
\(\langle\id_X,b\rangle\circ p\) pertany a la imatge directa, i pel
lema~\ref{lema:pertinenca-recobriment} també \(\langle\id_X,b\rangle\).
Recíprocament, si \(\langle\id_X,b\rangle\) pertany a
\((\id_X\times f)_*(s_\alpha)\), prenem el pullback de l'epimorfisme
\(s_\alpha\twoheadrightarrow\operatorname{Im}\) al llarg del testimoni;
dona un recobriment \(p\colon X'\twoheadrightarrow X\)
(lema~\ref{lema:estabilitat-epi-regular}) i un \(w\colon X'\to s_\alpha\),
i \(a:=\pi_A\circ s_\alpha\circ w\) compleix el que cal. Sobre punts
globals, un recobriment \(p\colon X'\twoheadrightarrow\mathbf1\) té un punt
(lema del buit i no-degeneració), i component-hi s'obté el testimoni
sobre \(\mathbf1\).

\emph{(3)} Usarem dos fets auxiliars. \emph{(a)} Per a una part
\(M\hookrightarrow Z\) i un element generalitzat \(z\colon Y\to Z\),
\(z\in_Y M^c\) si i només si la imatge inversa \(z^{-1}(M)\) és buida. En
efecte, \(\chi_{M^c}\circ z=\neg\circ\chi_M\circ z\)
(proposició~\ref{prop:caracteristica-complement}), que val \(1\circ{!}_Y\) si i
només si \(\chi_M\circ z=0\circ{!}_Y\) (componeu amb la involució \(\neg\)), és
a dir, si i només si \(\chi_{z^{-1}(M)}=\chi_{\bot_Y}\)
(proposicions~\ref{prop:caracteristica-imatge-inversa}
i~\ref{prop:caracteristica-buida}), que equival a \(z^{-1}(M)\equiv_Y\bot_Y\)
(proposició~\ref{prop:caracteristica-determina-part}). \emph{(b)} Si
\(p\colon T'\twoheadrightarrow T\) és un epimorfisme i \(T'\) és inicial,
\(T\) és inicial: \(0\circ{!}_T\circ p=1\circ{!}_T\circ p\) perquè \(T'\) és
inicial, d'on \(0\circ{!}_T=1\circ{!}_T\), i l'argument de la
proposició~\ref{prop:caracteristica-buida} dona una aplicació \(T\to\mathbf0\).

Escrivim \(K:=s_\alpha\hookrightarrow X\times A\) i
\(N:=(\id_X\times f)_*(K^c)\hookrightarrow X\times B\). Per~(1),
\(b\mathrel{\varepsilon}(\forall_f\circ\alpha)\) vol dir
\(\langle\id_X,b\rangle\in_X N^c\); i, pel
corol·lari~\ref{cor:pertinenca-universal},
\(a\mathrel{\varepsilon}(\alpha\circ x)\) vol dir \(\langle x,a\rangle\in_Y K\).

\(\Rightarrow\)) Suposem \(\langle\id_X,b\rangle\in_X N^c\), i siguin
\(x\colon Y\to X\), \(a\colon Y\to A\) amb \(f\circ a=b\circ x\). Sigui
\(t\colon T\hookrightarrow Y\) la imatge inversa \(\langle x,a\rangle^{-1}(K^c)\).
Aleshores \(\langle x,a\rangle\circ t\in K^c\) i, per tant,
\((\id_X\times f)\circ\langle x,a\rangle\circ t=\langle\id_X,b\rangle\circ x\circ t\)
pertany a la imatge directa \(N\). Però també pertany a \(N^c\), perquè
\(\langle\id_X,b\rangle\) hi pertany. Com que \(N\cap N^c\equiv\bot\), \(T\)
admet una aplicació cap a \(\mathbf0\) i és inicial (estrictesa). Per~(a),
\(\langle x,a\rangle\in_Y(K^c)^c\equiv K\), és a dir,
\(a\mathrel{\varepsilon}(\alpha\circ x)\).

\(\Leftarrow\)) Suposem la condició de l'enunciat, i sigui
\(t\colon T\hookrightarrow X\) la imatge inversa \(\langle\id_X,b\rangle^{-1}(N)\),
de manera que \(\langle t,b\circ t\rangle\in_T N\). Com a la prova de~(2),
el pullback de l'epimorfisme \(K^c\twoheadrightarrow N\) al llarg del
testimoni d'aquesta pertinença dona un epimorfisme
\(p\colon T'\twoheadrightarrow T\) i un element \(\langle x',a'\rangle\colon
T'\to X\times A\) de \(K^c\) amb \((\id_X\times f)\circ\langle x',a'\rangle
=\langle t,b\circ t\rangle\circ p\); és a dir, \(x'=t\circ p\) i
\(f\circ a'=b\circ x'\). La hipòtesi dona \(\langle x',a'\rangle\in_{T'}K\), i
com que també és a \(K^c\), \(T'\) és inicial. Per~(b), \(T\) és inicial, i
per~(a), \(\langle\id_X,b\rangle\in_X N^c\), és a dir,
\(b\mathrel{\varepsilon}(\forall_f\circ\alpha)\).
\end{prova}

\begin{proposicio}[Els quantificadors com a adjunts]\label{prop:quantificadors-adjuncio}
Sigui \(f\colon A\to B\). Per a tot estadi \(X\), tota
\(\alpha\colon X\to P(A)\) i tota \(\beta\colon X\to P(B)\), amb parts
associades \(s_\alpha\hookrightarrow X\times A\) i
\(s_\beta\hookrightarrow X\times B\), i escrivint \(\subseteq\) per a la
inclusió de parts:
\[
s_{\exists_f\alpha}\subseteq s_\beta
\ \Longleftrightarrow\
s_\alpha\subseteq s_{P(f)\beta},
\qquad
s_{P(f)\beta}\subseteq s_\alpha
\ \Longleftrightarrow\
s_\beta\subseteq s_{\forall_f\alpha} .
\]
És a dir, \(\exists_f\) és adjunt per l'esquerra i \(\forall_f\) adjunt
per la dreta de la imatge inversa interna \(P(f)\), a cada estadi.
\end{proposicio}

\begin{prova}
Per les proposicions~\ref{prop:imatge-inversa-interna}
i~\ref{prop:quantificadors-lectura}, \(s_{\exists_f\alpha}=(\id_X\times f)_*s_\alpha\)
i \(s_{P(f)\beta}=(\id_X\times f)^{-1}s_\beta\); la primera equivalència és,
doncs, l'adjunció imatge directa\,--\,imatge inversa del
capítol~\ref{cap:imatges} (proposició~\ref{prop:adjuncio-imatges}) per a
l'aplicació \(\id_X\times f\). La segona se'n dedueix per dualitat
booleana:
\(s_\beta\subseteq s_{\forall_f\alpha}=\bigl((\id\times f)_*(s_\alpha^c)\bigr)^c\)
equival, per l'antimonotonia del complement, a
\((\id\times f)_*(s_\alpha^c)\subseteq s_\beta^c\), i per la primera
adjunció a \(s_\alpha^c\subseteq(\id\times f)^{-1}(s_\beta^c)
=\bigl((\id\times f)^{-1}s_\beta\bigr)^c\) (la imatge inversa commuta amb
el complement, proposició~\ref{prop:imatge-inversa-complement}), que per
l'antimonotonia és \((\id\times f)^{-1}s_\beta\subseteq s_\alpha\).
\end{prova}

\begin{observacio}[Lectura clàssica]\label{obs:quantificadors-classica}
Per a \(f\colon A\to B\), la imatge inversa és
\(P(f)(\beta)=f^{-1}(\beta)=\{\,a:f(a)\in\beta\,\}\); el quantificador
existencial és la imatge directa
\(\exists_f(\alpha)=\{\,b:\exists a\,(f(a)=b\ \text{i}\ a\in\alpha)\,\}\);
i l'universal és
\(\forall_f(\alpha)=\{\,b:\forall a\,(f(a)=b\Rightarrow a\in\alpha)\,\}\).
Els noms «existencial» i «universal» no són metafòrics: quan
\(f={!}_A\colon A\to\mathbf1\), \(\exists_A\) diu si hi ha algun element de
\(A\) que pertany a la part, i \(\forall_A\) si tots hi pertanyen. La
lògica de predicats sobre \(A\), habitualment presa com a llenguatge
extern, apareix aquí com l'àlgebra de \(P(A)\) amb aquests adjunts. El
preu de llegir-ho sobre elements generalitzats és el recobriment de la
proposició~\ref{prop:quantificadors-lectura}: la mateixa frontera que a
la composició de relacions.
\end{observacio}

\section{\texorpdfstring{La inclusió com a ordre sobre \(P(A)\)}{La inclusió com a ordre sobre P(A)}}\label{sec:inclusio-ordre}

La inclusió de parts \(\subseteq_A\) del capítol~\ref{cap:parts} era una
relació del metallenguatge entre monomorfismes. Ara la realitzarem com una
\emph{part} de \(P(A)\times P(A)\), és a dir, com una relació interna
sobre el conjunt de parts, i mostrarem que és un ordre. L'eina és el
quantificador universal al llarg d'una projecció.

\begin{definicio}[Inclusió interna]\index{Inclusió interna}\label{def:inclusio-interna}
Sigui \(\Rightarrow:=\vee\circ(\neg\times\id_{\mathbf2})\colon
\mathbf2\times\mathbf2\to\mathbf2\) la implicació booleana, i sigui
\(I_A\hookrightarrow P(A)\times P(A)\times A\) la part classificada per
\(\Rightarrow\circ\langle\operatorname{ev}_1,\operatorname{ev}_2\rangle\).
La \emph{relació d'inclusió interna} sobre \(P(A)\) és la part
\[
{\subseteq_A}\;:=\;\forall_\pi(I_A)\;\hookrightarrow\;P(A)\times P(A),
\]
on \(\pi\colon P(A)\times P(A)\times A\to P(A)\times P(A)\) és la
projecció i \(\forall_\pi\) s'aplica al nom de la part \(I_A\) (és a dir,
\(\subseteq_A\) és la part de \(P(A)\times P(A)\) el nom de la qual és
\(\forall_\pi\ulcorner I_A\urcorner\)).
\end{definicio}

\begin{proposicio}[Lectura de la inclusió interna]\label{prop:inclusio-interna-lectura}
Siguin \(\alpha,\beta\colon X\to P(A)\), amb parts associades
\(s_\alpha,s_\beta\hookrightarrow X\times A\). Aleshores
\[
\begin{aligned}
\langle\alpha,\beta\rangle\in_X{\subseteq_A}
\ &\Longleftrightarrow\
s_\alpha\subseteq s_\beta\\
&\Longleftrightarrow\
\text{per a tot } x\colon Y\to X \text{ i tot } a\colon Y\to A,\\
&\qquad\qquad
a\mathrel{\varepsilon}(\alpha\circ x)\Rightarrow a\mathrel{\varepsilon}(\beta\circ x).
\end{aligned}
\]
En particular, sobre els noms de dues parts \(m,n\) de \(A\),
\(\langle\ulcorner m\urcorner,\ulcorner n\urcorner\rangle\in_{\mathbf1}{\subseteq_A}\)
si i només si \(m\subseteq_A n\).
\end{proposicio}

\begin{prova}
Per la lectura del quantificador universal
(proposició~\ref{prop:quantificadors-lectura}(3)) amb \(f=\pi\),
\(\langle\alpha,\beta\rangle\in_X\forall_\pi(I_A)\) vol dir que, per a tot
\(x\colon Y\to X\) i tot \(c\colon Y\to P(A)\times P(A)\times A\) amb
\(\pi\circ c=\langle\alpha,\beta\rangle\circ x\) ---és a dir, per a tot
\(c=\langle\alpha\circ x,\beta\circ x,a\rangle\) amb \(a\colon Y\to A\)---,
\(c\in_YI_A\), és a dir,
\(\Rightarrow\circ\langle u,v\rangle=1\circ{!}_Y\) amb
\(u=\operatorname{ev}\circ\langle\alpha\circ x,a\rangle\) i
\(v=\operatorname{ev}\circ\langle\beta\circ x,a\rangle\). Ara bé,
\(\Rightarrow\circ\langle u,v\rangle=\neg u\vee v\) és la característica
de la part \(T^c\cup U\) de \(Y\), on \(T,U\) són les parts classificades
per \(u,v\); i \(T^c\cup U\) és total si i només si \(T\subseteq U\)
(lleis booleanes, teorema~\ref{teo:lleis-booleanes}). Per la
proposició~\ref{prop:pertinenca-interna-lectura}, \(T\) i \(U\) són les
imatges inverses \(\langle x,a\rangle^{-1}s_\alpha\) i
\(\langle x,a\rangle^{-1}s_\beta\). Que la condició valgui per a tots els
\(x,a\) és, prenent \(Y=s_\alpha\) i \(\langle x,a\rangle\) la seva
inclusió, exactament \(s_\alpha\subseteq s_\beta\); i recíprocament,
\(s_\alpha\subseteq s_\beta\) la implica per monotonia de la imatge
inversa. Això dona les dues equivalències (la tercera és la segona,
llegida amb el corol·lari~\ref{cor:pertinenca-universal}). Per als noms,
\(X=\mathbf1\), i \(s_{\ulcorner m\urcorner}\) és \(m\) vista dins
\(\mathbf1\times A\cong A\).
\end{prova}

\begin{teorema}[La inclusió interna és un ordre parcial]\label{teo:inclusio-ordre}
La relació \(\subseteq_A\) sobre \(P(A)\) és reflexiva, antisimètrica i
transitiva, en el sentit intern següent, per a tots
\(\alpha,\beta,\gamma\colon X\to P(A)\):
\begin{enumerate}
\item \(\langle\alpha,\alpha\rangle\in_X{\subseteq_A}\);
\item si \(\langle\alpha,\beta\rangle\in_X{\subseteq_A}\) i
\(\langle\beta,\alpha\rangle\in_X{\subseteq_A}\), aleshores
\(\alpha=\beta\);
\item si \(\langle\alpha,\beta\rangle\in_X{\subseteq_A}\) i
\(\langle\beta,\gamma\rangle\in_X{\subseteq_A}\), aleshores
\(\langle\alpha,\gamma\rangle\in_X{\subseteq_A}\).
\end{enumerate}
A més, la intersecció i la unió internes són l'ínfim i el suprem respecte
d'aquest ordre: \(\langle\gamma,\alpha\cap\beta\rangle\in_X{\subseteq_A}\)
si i només si \(\langle\gamma,\alpha\rangle\in_X{\subseteq_A}\) i
\(\langle\gamma,\beta\rangle\in_X{\subseteq_A}\), i dualment per a
\(\cup\).
\end{teorema}

\begin{prova}
Per la proposició~\ref{prop:inclusio-interna-lectura}, tot es redueix a la
inclusió de parts de \(X\times A\): \(s_\alpha\subseteq s_\alpha\); si
\(s_\alpha\subseteq s_\beta\) i \(s_\beta\subseteq s_\alpha\), les parts
són equivalents i tenen la mateixa característica, de manera que
\(\alpha=\beta\) per la unicitat de la transposada; i la transitivitat de
la inclusió de parts és immediata de la
proposició~\ref{prop:inclusio-pertinenca}. Per a
l'ínfim, \(s_{\alpha\cap\beta}=s_\alpha\cap s_\beta\) per la
proposició~\ref{prop:booleanes-puntuals} (les característiques
coincideixen), i una part és inclosa en una intersecció si i només si ho
és en cadascun dels factors (proposició~\ref{prop:pertinenca-interseccio});
dualment, \(s_{\alpha\cup\beta}=s_\alpha\cup s_\beta\) i la unió és el
suprem per les lleis booleanes (proposició~\ref{prop:inclusio-booleana}).
\end{prova}

\begin{observacio}[Lectura clàssica]\label{obs:inclusio-ordre-classica}
El teorema diu que \((P(A),\subseteq)\) és un reticle booleà: un conjunt
parcialment ordenat per la inclusió, amb la intersecció com a ínfim, la
unió com a suprem, la part buida i la total com a mínim i màxim, i el
complement com a operació de complementació. És l'estructura de reticle de
potència de la teoria clàssica, ara construïda com una relació interna
sobre l'objecte \(P(A)\) en lloc de descrita externament sobre la
col·lecció de subconjunts.
\end{observacio}

\begin{resumcap}
\apartat{Idees centrals}
\begin{enumerate}
  \item Tres nivells: la part \(m\colon S\hookrightarrow A\) (una fletxa), el seu nom \(\ulcorner m\urcorner\colon\mathbf1\to P(A)\) (un punt) i el conjunt \(P(A)\). La pertinença interna \(\in_A\hookrightarrow A\times P(A)\) és la part classificada per l'avaluació; sobre noms recupera la pertinença generalitzada.
  \item Una aplicació \(X\to P(A)\) és el mateix que una part de \(X\times A\): una família de parts d'\(A\) indexada per \(X\). Tot el capítol llegeix les construccions sobre \(P(A)\) a través d'aquesta correspondència.
  \item Les operacions booleanes internes es defineixen per transposició i es calculen component avaluacions; la lectura amb \enquote{i} val a tots els estadis, les lectures amb \enquote{o} i \enquote{no} només sobre punts. Sobre noms recuperen les operacions de parts, i les lleis booleanes es transporten.
  \item \(P(f)\) és la imatge inversa interna; \(\exists_f\) és la transposada de la imatge directa de la pertinença al llarg de \(\id\times f\), i \(\forall_f:=\neg\exists_f\neg\). La imatge directa commuta amb el canvi de base (Beck--Chevalley), i això dona la lectura de \(\exists_f\) sobre elements generalitzats (amb recobriment) i l'adjunció \(\exists_f\dashv P(f)\dashv\forall_f\) a cada estadi.
  \item La inclusió interna és \(\forall_\pi\) de la implicació puntual; es llegeix com la inclusió de les parts associades, i és un ordre amb \(\cap\) i \(\cup\) com a ínfim i suprem.
  \item Cap axioma nou: tot surt dels exponencials, el classificador i la factorització imatge.
\end{enumerate}
\apartat{Errors típics} Identificar la part \(m\) amb el seu nom \(\ulcorner m\urcorner\) o amb el conjunt \(S\); llegir \(a\mathrel{\varepsilon}(\alpha\cup\beta)\) com \enquote{\(a\mathrel{\varepsilon}\alpha\) o \(a\mathrel{\varepsilon}\beta\)} per a un element generalitzat; creure que \(\exists_f\) \enquote{existeix} per la seva lectura puntual, quan és una aplicació que cal construir amb la imatge directa; esperar el testimoni de \(\exists_f\) en el mateix estadi sense elecció; definir la inclusió interna amb \(\forall_A\) (que quantifica cap a \(\mathbf2\)) en lloc de \(\forall_\pi\); i oblidar que \(\forall_f=\neg\exists_f\neg\) depèn del caràcter booleà de la teoria.
\apartat{Lectura recomanada} \cite[cap.~7--8]{lawvere-rosebrugh} per al conjunt de parts, la pertinença interna, els quantificadors com a adjunts i el reticle de parts; \cite[article~VI, sessions~32--33]{lawvere-schanuel} per a una presentació elemental de l'objecte de parts; \cite[cap.~9]{awodey} per a les adjuncions en general, amb els quantificadors com a exemple; \cite[cap.~IV]{maclane-moerdijk} per als objectes de parts en un topos i la condició de Beck--Chevalley; i \cite[§§5--9]{halmos} per al conjunt de parts i les operacions amb elements.
\end{resumcap}
